\documentclass[12pt]{article}

\usepackage{amsmath}
\usepackage{amssymb}
\usepackage{physics}
\usepackage{hyperref}

\title{Affine Independence and Spans}
\author{}
\date{}

\begin{document}

\maketitle

\section{Affine Subspaces Spanned by Points}

\begin{definition}
Let \(\ell\geq 1\). The points
\[
x_1,\ldots,x_\ell\in \mathcal A
\]
are called \emph{affinely independent} if there exists no affine
subspace of \(\mathcal A\) of dimension at most \(\ell-2\) containing all of them.
Otherwise, the points are called \emph{affinely dependent}.
\end{definition}

For example, suppose that the ambient affine space \(\mathcal A\) has dimension
\(n=3\).

\begin{itemize}
    \item If \(\ell=1\), a single point is always affinely independent.

    \item If \(\ell=2\), two points are affinely independent if and only if
    they are distinct.

    \item If \(\ell=3\), three points are affinely independent if and only if
    they do not lie on the same affine line.

    \item If \(\ell=4\), four points are affinely independent if and only if
    they do not lie in the same affine plane.

    \item If \(\ell\geq 5\), the points are always affinely dependent.
\end{itemize}

Indeed, when \(n=3\) and \(\ell\geq5\), we have
\[
\ell-2\geq3.
\]
Since all the points lie in \(\mathcal A\), which is itself a
three-dimensional affine space, they are affinely dependent.

\begin{proposition}
There exist \(\ell\) affinely independent points of \(\mathcal A\) whenever
\[
1\leq \ell\leq n+1,
\]
where \(n=\dim \mathcal A=\dim V\).
\end{proposition}

\begin{proof}
Choose a point \(x\in \mathcal A\). Since \(\dim V=n\), we may choose a basis
\[
\mathbf v_1,\ldots,\mathbf v_n
\]
of \(V\). Define the points
\[
x_i=\mathbf v_i(x),
\qquad i=1,\ldots,n.
\]

For \(\ell=1\), the set \(\{x\}\) is contained in the zero-dimensional
affine subspace
\[
S(x,\{\mathbf 0\})=\{x\},
\]
and \(x\) is affinely independent.

Now let
\[
2\leq \ell\leq n+1
\]
and define
\[
U=\operatorname{span}
\{\mathbf v_1,\ldots,\mathbf v_{\ell-1}\}.
\]
Since the vectors
\[
\mathbf v_1,\ldots,\mathbf v_{\ell-1}
\]
belong to a basis of \(V\), they are linearly independent. Therefore,
\[
\dim U=\ell-1.
\]

The points
\[
x,x_1,\ldots,x_{\ell-1}
\]
belong to the affine subspace \(S(x,U)\); that is,
\[
\{x,x_1,\ldots,x_{\ell-1}\}\subseteq S(x,U).
\]

Suppose that these points also belong to an arbitrary affine subspace
containing \(x\). Since \(x\) belongs to this affine subspace, we may
write it as \(S(x,W)\) for some linear subspace \(W\leq V\). Thus,
\[
\{x,x_1,\ldots,x_{\ell-1}\}
\subseteq S(x,U)\cap S(x,W).
\]

Because the intersection contains \(x\), it is nonempty, and hence
\[
S(x,U)\cap S(x,W)=S(x,U\cap W).
\]
Consequently,
\[
x_i\in S(x,U\cap W),
\qquad i=1,\ldots,\ell-1.
\]

Since
\[
x_i=\mathbf v_i(x),
\]
the uniqueness of the vector mapping \(x\) to \(x_i\) implies that
\[
\mathbf v_i\in U\cap W,
\qquad i=1,\ldots,\ell-1.
\]
Therefore,
\[
U
=\operatorname{span}
\{\mathbf v_1,\ldots,\mathbf v_{\ell-1}\}
\subseteq U\cap W.
\]
Since \(U\cap W\subseteq U\), we conclude that
\[
U\cap W=U.
\]
In particular,
\[
U\subseteq W,
\]
and hence
\[
\dim W\geq\dim U=\ell-1.
\]

Thus every affine subspace containing
\[
x,x_1,\ldots,x_{\ell-1}
\]
has dimension at least \(\ell-1\). Therefore, these \(\ell\) points cannot
all belong to an affine subspace of dimension at most \(\ell-2\), and
they are affinely independent.
\end{proof}

\begin{proposition}
Let
\[
1\leq \ell\leq n+1.
\]
If the points
\[
x_1,\ldots,x_\ell\in \mathcal A
\]
are affinely independent, then there exists one and only one
\((\ell-1)\)-dimensional affine subspace of \(\mathcal A\) containing them.
\end{proposition}

\begin{proof}
If \(\ell=1\), the unique zero-dimensional affine subspace containing
\(x_1\) is
\[
S(x_1,\{\mathbf 0\})=\{x_1\}.
\]

Let now \(\ell\geq2\), and consider the \(\ell-1\) vectors
\[
\mathbf d_i=\overrightarrow{x_1x_{i+1}},
\qquad i=1,\ldots,\ell-1.
\]

We first show that these vectors are linearly independent. Suppose,
on the contrary, that they are linearly dependent, and let
\[
U=\operatorname{span}
\{\mathbf d_1,\ldots,\mathbf d_{\ell-1}\}.
\]
Their linear dependence would imply that
\[
\dim U\leq \ell-2.
\]

By construction,
\[
x_{i+1}=\mathbf d_i(x_1),
\qquad i=1,\ldots,\ell-1.
\]
Therefore,
\[
x_1,\ldots,x_\ell\in S(x_1,U).
\]
The points would consequently lie in an affine subspace of dimension
at most \(\ell-2\), contradicting their affine independence. Hence,
\[
\mathbf d_1,\ldots,\mathbf d_{\ell-1}
\]
are linearly independent.

It follows that
\[
\dim U=\ell-1,
\]
and therefore
\[
S=S(x_1,U)
\]
is a \((\ell-1)\)-dimensional affine subspace containing
\(x_1,\ldots,x_\ell\). This proves existence.

We now prove uniqueness. Suppose that \(S'\) is another
\((\ell-1)\)-dimensional affine subspace containing the same points.
Write
\[
S'=S(y,W),
\qquad
\dim W=\ell-1.
\]
Since \(x_1\in S'\), we may take \(x_1\) as its reference point, and
hence
\[
S'=S(x_1,W).
\]

All the points \(x_1,\ldots,x_\ell\) belong to both \(S(x_1,U)\) and
\(S(x_1,W)\). Thus,
\[
x_1,\ldots,x_\ell
\in
S(x_1,U)\cap S(x_1,W)
=
S(x_1,U\cap W).
\]

For every \(i=1,\ldots,\ell-1\), we therefore have
\[
\mathbf d_i
=
\overrightarrow{x_1x_{i+1}}
\in U\cap W.
\]
Since the vectors \(\mathbf d_1,\ldots,\mathbf d_{\ell-1}\) span \(U\),
it follows that
\[
U\subseteq U\cap W.
\]
The reverse inclusion is automatic, so
\[
U\cap W=U.
\]
Consequently,
\[
U\subseteq W.
\]

But
\[
\dim U=\dim W=\ell-1,
\]
and therefore
\[
U=W.
\]
It follows that
\[
S'=S(x_1,W)=S(x_1,U)=S.
\]
Hence the \((\ell-1)\)-dimensional affine subspace containing
\(x_1,\ldots,x_\ell\) is unique.
\end{proof}

\end{document}
